\section{Physical Measurement System}
\label{sec:system}

\subsection{Instrumentation, Timing, and Calibration}
The instrument comprises a two-view camera with an endoscopic field of view,
a camera-mounted WitMotion USB IMU, a GP050 board, and a second WitMotion BLE
IMU fixed to the board. The USB unit measures the camera rigid body; the BLE
unit is not a camera input and only verifies that the reference board is
stationary. Table~\ref{tab:instrument} lists the parameters entering the
measurement model together with their provenance.

\begin{table}[t]
\centering
\caption{Instrument parameters and measured provenance.}
\label{tab:instrument}
\footnotesize
\setlength{\tabcolsep}{4pt}
\begin{tabular}{lll}
\toprule
Quantity & Value & Source \\
\midrule
Board inner corners & $11\times8$ & GP050 specification \\
Square pitch & 3~mm & Board specification \\
Left intrinsic RMS & 0.243~px & 80 calibration images \\
Right intrinsic RMS & 0.284~px & 80 calibration images \\
Nominal baseline & 46.5~mm & Mechanical measurement \\
USB IMU rate & 200.0~Hz & Session timestamps \\
BLE IMU rate & 187.0--188.0~Hz & Session timestamps \\
Camera delay & 14--90~ms & Motion correlation \\
\bottomrule
\end{tabular}
\end{table}

\paragraph{Concurrent acquisition and timing.}
Camera grabbing, USB decoding, BLE notification processing, board detection,
and disk writing execute in separate threads pinned to distinct logical CPU
cores. The accepted recordings have a median USB interval near 5.0~ms and no
gap above 20~ms. Camera times are host-clock arrival times and include
exposure and USB delivery. A per-recording delay is estimated by maximizing
correlation between image-difference and gyroscope magnitudes. Interpolation
is used when an IMU value is required on a regular 200-Hz grid; frame metadata
records the latest sample at or before the delay-corrected frame time.

\paragraph{Calibration.}
Both cameras are calibrated at $1280\times720$ using Zhang's
method~\cite{zhang2000flexible}; matrices are scaled for other resolutions
while distortion is retained. IMU-to-camera rotation is estimated in a
dedicated three-axis recording by Wahba/Kabsch fitting on one half of
one-second increments and evaluated on the other half
(Section~\ref{sec:imu}). A fixed-intrinsic stereo fit on quasi-static frames
of the intrinsic recording gives a 46.3-mm baseline (nominal 46.5~mm) and a
$10.2^\circ$ roll between the two cameras, with 0.58-px reprojection RMS. The
cameras are not hardware-synchronized, and the fit does not meet our held-out
acceptance limits; right images therefore support rendering and acquisition
checks but not the reference pose.

\subsection{Real Dataset and Audit}
The split contains nine complete trajectory recordings totaling 871.5~s:
six training recordings (495.5~s), two validation recordings (260.8~s), and
one held-out test recording (115.2~s). Acceptance requires at least 40~s,
USB rate $\geq150$~Hz with at most five gaps above 50~ms, BLE rate
$\geq120$~Hz, and left images for at least 85\% of frames. No window from one
recording appears in another split. At the selected 3-s horizon, train,
validation, and test contain 6754, 1436, and 647 pairs, respectively. Four
rigidity recordings (342~s) that passed the mount test also carry chessboard
poses; they contain nodding, yawing, rolling, and still phases with little
translation and are used as training-only data in the cross-validation
protocol. Two rigidity recordings made before the mount was secured fail the
test and are excluded from every split.

Fig.~\ref{fig:dataset} shows the corpus qualitatively. The board occupies a
small part of the endoscopic field of view at the working distance, which is
why the corner-noise propagation of Section~\ref{sec:reference} matters;
panels (e) and (f) show the 200-Hz stream together with the reference speed
that defines the stillness labels, and panel (g) shows that the measurand
distribution differs between splits---the held-out test recording is the fastest
of the nine, with a mean 3-s displacement of 50~mm against 28~mm for the
training recordings, so its zero-motion baseline is correspondingly larger.

\begin{figure}[tb]
\centering
\includegraphics[width=\columnwidth]{dataset.pdf}
\caption{The corpus. (a)--(c) Real acquisitions, cropped to the board with
detected corners in green; (d) synthetic renders. (e), (f) One inertial
interval, with reference-confirmed pauses and the two stillness thresholds.
(g) The held-out recording has a larger 3-s displacement than the training set.}
\label{fig:dataset}
\end{figure}
